% SPDX-License-Identifier: Apache-2.0
\section{A Frame's Length, Found Again}
\label{sec:x-framelen}

\code{FrameIn} (\S\ref{sec:x-ethdma}) must know a frame's length
before the store engine starts, since the engine is told how many
bytes to write when it begins. \code{EthRx} knows it and says so. On
the board it cannot, and the reason is the clock.

The board's receiver is in the top, on the clock the PHY recovers from
the line, and only the frame's bytes cross to the core's clock. The
length cannot follow them on a wire of its own: it would arrive on the
wrong clock, out of step with the frame it describes. Nor can it ride
inside the stream, because the remote peripheral reads that same
stream, through the sharing unit of \S\ref{sec:x-ethshare}, and would
take the length for a byte of its own frame.

So the length is found again after the crossing. \code{FrameLen} takes
a frame a byte at a time, holds all of it, and then offers it back with
its count on \code{len}, which is exactly what \code{EthRx} offers. A
unit joined to \code{FrameIn} in its place cannot tell the difference.

\lstinputlisting[rangebeginprefix=//\ begin\{,rangebeginsuffix=\},%
rangeendprefix=//\ end\{,rangeendsuffix=\},includerangemarker=false,%
linerange=len-len,caption={\code{lib/parts/src/ethdma.rs}: a frame held
whole and handed on with its count.},%
label={lst:x-framelen}]{lib/parts/src/ethdma.rs}

\subsection{What it costs}

A frame of memory and a frame of latency. The frame was already held
whole once, before the crossing, by the receiver; it is held whole
again here. That is the price of the length not crossing with the
bytes, and it is paid in block RAM rather than in a second clock
domain's worth of logic.

\subsection{What the run checks}

The run is \code{FrameOut}, then \code{FrameLen}, then \code{FrameIn},
with two frames of 101 and 67 bytes, the same two \code{ex\_ethdma}
sends through the real MAC, so the two runs can be read side by side.
Every word is checked where it landed, and the length \code{FrameIn}
was told is checked against the frame's.

Made to report every length one short, \code{FrameLen} fails the run at
once: the first frame's last word comes back without its last byte, and
the second frame arrives shifted by the byte the first one left behind.
A length that is wrong by one is the fault this unit exists to prevent,
and the run is built to see it.

\input{framelen_timing.tex}

\lstinputlisting[style=ascii,caption={What \code{ex\_framelen} printed
when the build ran it.},%
label={lst:x-framelen-out}]{out_framelen.txt}
